\documentclass[11pt,a4paper]{article}

\usepackage[margin=2cm]{geometry}
\usepackage[hidelinks]{hyperref}
\usepackage[T1]{fontenc}
\usepackage[utf8]{inputenc}

\setlength{\parindent}{0pt}
\pagenumbering{gobble}

\begin{document}

\begin{center}
{\Large \textbf{MIGUEL ANGEL ANAY}}\\
Lima, Perú\\
Software Engineer | AI Architect | Backend \& Cloud Engineer\\
\href{mailto:contacto@miguel-anay.nom.pe}{contacto@miguel-anay.nom.pe} \;|\;
\href{https://www.linkedin.com/in/miguel-anay/}{Linkedin} \;|\;
\href{https://portfolio.miguel-anay.nom.pe}{Portafolio}
\end{center}

\vspace{0.3cm}

\section*{Perfil Profesional}

Ingeniero de Software con más de \textbf{12 años de experiencia} en desarrollo backend, arquitectura de software y soluciones empresariales. Especializado en \textbf{arquitecturas de Inteligencia Artificial, LLMs, RAG, Cloud Computing (Azure y AWS)} y diseño de sistemas escalables. Experiencia liderando equipos técnicos, definiendo estándares de desarrollo, buenas prácticas y soluciones orientadas a negocio.

\section*{Experiencia Profesional}

\textbf{Encora} \hfill \textit{Nov 2020 – Abr 2025}\\
\textit{Software Engineer III}

\begin{itemize}
  \item Desarrollo y mantenimiento de plataformas empresariales para Profuturo (Agencia Virtual, Clave Web, Disfruta).
  \item Diseño de arquitecturas backend escalables utilizando C\#.NET, Node.js y Python.
  \item Implementación de principios SOLID, arquitectura limpia y patrones de diseño.
  \item Integración de APIs, servicios cloud y controles de seguridad.
  \item Liderazgo técnico, code reviews y soporte a equipos de desarrollo.
\end{itemize}

\textbf{Prosegur} \hfill \textit{2018 – 2020}\\
\textit{Software Engineer}

\begin{itemize}
  \item Desarrollo de sistemas críticos como Control de Tiempos Operativos (CTO).
  \item Implementación de soluciones backend con C\#.NET, SQL Server y MongoDB.
  \item Desarrollo de plataformas Web Audit, Web PER y aplicaciones móviles de supervisión.
  \item Modelamiento de datos y optimización de consultas.
\end{itemize}

\textbf{Medios Industriales} \hfill \textit{2017 – 2018}\\
\textit{Software Developer}

\begin{itemize}
  \item Desarrollo de sistemas de control de tiempos con C\#.NET y MySQL.
\end{itemize}

\section*{Inteligencia Artificial y Arquitecturas Cloud}

\begin{itemize}
  \item Diseño e implementación de soluciones basadas en LLMs y RAG (Retrieval-Augmented Generation).
  \item Experiencia en chunking, embeddings y bases de datos vectoriales.
  \item Integración de IA con lógica de negocio y clasificación de intención.
  \item Uso de Azure OpenAI Service, Azure AI Search y AWS Bedrock.
  \item Diseño de arquitecturas orientadas a agentes y flujos inteligentes.
\end{itemize}

\section*{Tecnologías}

\textbf{Lenguajes:} Python, C\#.NET, TypeScript, JavaScript\\
\textbf{Backend:} FastAPI, Flask, Django, Node.js\\
\textbf{IA:} Azure OpenAI, AWS Bedrock, LangChain, LangGraph, RAG\\
\textbf{Cloud:} Azure, AWS, GCP\\
\textbf{DevOps:} Docker, Kubernetes, CI/CD, Terraform\\
\textbf{Bases de Datos:} SQL Server, MySQL, MongoDB, Bases vectoriales\\
\textbf{Arquitectura:} Clean Architecture, MVC, DDD, Event-Driven\\
\textbf{Versionamiento:} Git, GitHub

\section*{Educación}

\textbf{Universidad Nacional de Ingeniería (UNI)}\\
Ingeniero Industrial \hfill \textit{2004 – 2012}\\
Colegio de Ingenieros del Perú (CIP): 301084

\vspace{0.2cm}

\textbf{Universidad Privada del Norte}\\
Diplomado en Gerencia de TI con Mención en Disrupción Tecnológica\\
\textit{Enero 2025 – Septiembre 2025}

\section*{Especialización Técnica y Cursos}

\textbf{Arquitectura Hexagonal C\#.NET}\\
Technical Track \hfill \textit{Octubre 2025 – 8.5 horas}

\textbf{Advanced Data Engineer }\\
DMC Instituto (Diplomado) \hfill \textit{Septiembre 2025 – Actual (144 horas)}

\textbf{Coursera Generative AI Training Path}\\
AI Technical Track Certification \hfill \textit{Marzo 2025 – 60 horas}

\textbf{Azure Fundamentals AZ-900}\\
Udemy / Microsoft Azure Certification \hfill \textit{Noviembre 2024 – 9 horas}

\textbf{Administración de SQL Server}\\
Sistemas UNI – Módulo de Base de Datos \hfill \textit{Junio 2015 – 96 horas}

\textbf{Experto en Business Intelligence}\\
Institución de Formación Especializada \hfill \textit{Febrero 2016 – 48 horas}

\textbf{Data Science – Python Data Science}\\
DSRP Certification Program \hfill \textit{Octubre 2025 – 20 horas}

\textbf{Visual Basic .NET}\\
Institución de Formación Técnica \hfill \textit{Septiembre 2011 – 48 horas}

\section*{Habilidades Clave}

Arquitectura de Software, Inteligencia Artificial, Liderazgo Técnico, Cloud Computing, Backend, DevOps, Modelamiento de Datos, Seguridad de Software, Optimización de Sistemas.

\end{document}
